\pagestyle{empty}
\tight
\begin{tblr}{width=\textwidth,colspec={|Q[c,m,1.9cm]|X[l,m]|},hlines,vlines,hspan=minimal,row{1}={bg=headblue},rowsep=1pt,colsep=3pt}
\SetCell{c,m}\hfil\logo\hfil & \textbf{POLITEKNIK NEGERI MALANG}\par\textbf{JURUSAN TEKNOLOGI INFORMASI}\par\textbf{PROGRAM STUDI: D4 TEKNIK INFORMATIKA} \\
\SetCell[c=2]{c} \textbf{RENCANA PEMBELAJARAN SEMESTER (RPS)} & \\
\end{tblr}
\begin{longtblr}{width=\textwidth,colspec={|Q[l,m,1.9cm]|X[0.85,l,m]|X[1.45,l,m]|X[0.75,l,m]|X[1.15,l,m]|},hlines,vlines,hspan=minimal,rowsep=1pt,colsep=2pt,row{1,3}={bg=headgray,font=\bfseries}}
MATA KULIAH & KODE & BOBOT (SKS)/jam & SEMESTER & TGL. PENYUSUNAN \\
Komputasi Hijau & RTI256108 & 2 SKS/4 jam & 6 & 1 Juli 2025 \\
OTORISASI & Dosen Pengembang RPS & Koordinator RMK & \SetCell[c=2]{l} Ka PRODI &  \\
 & Farid Angga Pribadi, S.Kom., M.Kom. &  & \SetCell[c=2]{l}{\vspace{8mm}\par Dr. Ely Setyo Astuti, S.T., M.T.} &  \\
\SetCell[r=9]{l} Capaian Pembelajaran (CP) & \SetCell[c=4]{l,bg=headgray,font=\bfseries} Capaian Pembelajaran Lulusan Program Studi (CPL-Prodi) &  &  &  \\
 & CPL01 & \SetCell[c=3]{l} Mampu menerapkan nilai ketakwaan, kesadaran hukum, kedisiplinan, profesionalisme, etika profesi, dan pembelajaran sepanjang hayat, serta tanggap terhadap isu sosial dan perkembangan teknologi. &  &  \\
 & CPL05 & \SetCell[c=3]{l} Mampu menghasilkan solusi teknologi komputasi dan infrastruktur digital yang sesuai dengan kebutuhan organisasi dan pengguna. &  &  \\
 & \SetCell[c=4]{l,bg=headgray,font=\bfseries} Capaian Pembelajaran mata kuliah (CPL-MK) &  &  &  \\
 & CPMK0101 & \SetCell[c=3]{l} Mahasiswa mampu menerapkan etika profesi TI dan tanggung jawab sosial termasuk aspek keberlanjutan lingkungan. &  &  \\
 & CPMK0509 & \SetCell[c=3]{l} Mahasiswa mampu merancang dan mengimplementasikan layanan komputasi virtual yang efisien dan berkelanjutan dari perspektif lingkungan. &  &  \\
 & \SetCell[c=4]{l,bg=headgray,font=\bfseries} Kemampuan akhir tiap tahapan belajar (Sub-CPMK) &  &  &  \\
 & SCPMK0101-05201 & \SetCell[c=3]{l} Mahasiswa mampu menganalisis dampak lingkungan dari infrastruktur TI dan merekomendasikan strategi komputasi hijau berbasis etika profesi. &  &  \\
 & SCPMK0509-05202 & \SetCell[c=3]{l} Mahasiswa mampu merancang dan mengevaluasi solusi infrastruktur TI hemat energi menggunakan virtualisasi, optimasi server, dan renewable energy. &  &  \\
\SetCell[r=2]{l} Penilaian & \SetCell[c=4]{l} *Nilai CPMK didapat dari agregasi nilai SUB-CPMK yang dibebankan &  &  &  \\
 & \SetCell[c=4]{l}{\tiny\begin{tblr}{width=\linewidth,colspec={|Q[c,m,0.1\linewidth]|Q[c,m,0.16\linewidth]|X[l,m]|X[l,m]|X[l,m]|X[l,m]|X[l,m]|X[l,m]|Q[c,m,0.08\linewidth]|},hlines,vlines,hspan=minimal,rowsep=1pt,colsep=0.7pt,row{1,2}={bg=headgray,font=\bfseries}}
Kode CPMK & Kode SCPMK & \SetCell[c=6]{c} Bobot per Bentuk Penilaian & & & & & & Total Bobot \\
 & & Tugas & Kuis 1 & UTS & Kuis 2 & PBL & UAS & \\
CPMK0101 & SCPMK0101-05201 & 15\%: dampak lingkungan TI & 5\%: konsep green IT & 8\%: carbon footprint & 0\% & 8\%: audit green IT & 0\% & 36\% \\
CPMK0509 & SCPMK0509-05202 & 15\%: infrastruktur hemat energi & 0\% & 7\%: efisiensi energi & 5\%: virtualisasi green & 14\%: audit green IT & 23\%: desain infrastruktur & 64\% \\
\SetCell[c=8]{r}\textbf{Total Per Penilaian} & & & & & & & & \textbf{100\%} \\
\end{tblr}} &  &  &  \\
Diskripsi Singkat Mata Kuliah & \SetCell[c=4]{l} Komputasi Hijau membangun kesadaran dan kemampuan merancang infrastruktur TI ramah lingkungan melalui efisiensi energi, virtualisasi, pengelolaan e-waste, dan penerapan prinsip keberlanjutan dalam pengembangan sistem TI. &  &  &  \\
Bahan Kajian / Materi Pembelajaran & \SetCell[c=4]{l}\begin{enumerate}\item Konsep komputasi hijau, carbon footprint, dan regulasi keberlanjutan TI\item Efisiensi energi perangkat keras, PUE, dan data center hijau\item Virtualisasi, green software, cloud computing, dan renewable energy\item E-waste, green IT governance, dan studi kasus sustainabilitas\end{enumerate} &  &  &  \\
\SetCell[r=2]{l} Pustaka & \SetCell[c=4]{l} \textbf{Utama:}\begin{enumerate}\item Murugesan, S. 2008. Harnessing Green IT. IEEE IT Professional.\item Belady, C. 2010. Green Grid Metrics. The Green Grid.\item Gupta, M. 2011. Green IT Strategy and Implementation. Springer.\end{enumerate} &  &  &  \\
 & \SetCell[c=4]{l} \textbf{Pendukung:} Materi LMS, dokumentasi standar ISO 14001, laporan The Green Grid, dan jobsheet praktikum. &  &  &  \\
Media Pembelajaran & \SetCell[c=2]{l} \textbf{Software:} Tools pengukuran energi, simulator data center, LMS. &  & \SetCell[c=2]{l} \textbf{Hardware:} Komputer/laptop, proyektor. &  \\
Nama Dosen Pengampu & \SetCell[c=4]{l} Tim Pengajar Program Studi D4 TI &  &  &  \\
Matakuliah Syarat & \SetCell[c=4]{l} - &  &  &  \\
\end{longtblr}
\begin{longtblr}{width=\textwidth,colspec={|Q[c,m,0.06\textwidth]|X[1.35,l,m]|X[1.08,l,m]|X[1.16,l,m]|X[0.7,l,m]|X[1.24,l,m]|X[1.06,l,m]|X[0.98,l,m]|Q[c,m,0.05\textwidth]|},hlines,vlines,hspan=minimal,rowhead=2,row{1,2}={bg=headgreen,font=\bfseries},rowsep=1pt,colsep=1.5pt}
Mgg.\par Ke & Kemampuan Akhir Yang Direncanakan (Sub-CP-MK) & Bahan kajian (Materi Pembelajaran) & Bentuk dan Metode Pembelajaran & Estimasi Waktu & Pengalaman Belajar Mahasiswa & Kriteria \& Bentuk Penilaian & Indikator Penilaian & Bobot Penilaian (\%) \\
(1) & (2) & (3) & (4) & (5) & (6) & (7) & (8) & (9) \\
1 & SCPMK0101-05201 & Pengantar komputasi hijau: definisi dan isu lingkungan TI. & Modalitas: Blended Learning. Bentuk: Luring/Daring. Metode: Case Method, studi kasus, diskusi lingkungan, dan peer review. & 1 x 2 x 50' tatap muka; 1 x 2 x 50' tugas/praktik mandiri. & Mahasiswa mengerjakan aktivitas terarah untuk pengantar komputasi hijau dan menunjukkan evidence hasil belajar. & Bentuk penilaian: Pengantar Komputasi Hijau. Mengacu rubrik RTM. & Ketepatan konsep; kualitas analisis; validasi dan dokumentasi. & 2 \\
2 & SCPMK0101-05201 & Carbon footprint dan pengukuran emisi TI. & Modalitas: Blended Learning. Bentuk: Luring/Daring. Metode: Case Method, studi kasus, diskusi lingkungan, dan peer review. & 1 x 2 x 50' tatap muka; 1 x 2 x 50' tugas/praktik mandiri. & Mahasiswa mengerjakan aktivitas terarah untuk carbon footprint dan pengukuran emisi TI dan menunjukkan evidence hasil belajar. & Bentuk penilaian: Carbon Footprint dan Pengukuran Emisi TI. Mengacu rubrik RTM. & Ketepatan konsep; kualitas analisis; validasi dan dokumentasi. & 2 \\
3 & SCPMK0101-05201 & Standar keberlanjutan TI: ISO 14001, Green Grid. & Modalitas: Blended Learning. Bentuk: Luring/Daring. Metode: Case Method, studi kasus, diskusi lingkungan, dan peer review. & 1 x 2 x 50' tatap muka; 1 x 2 x 50' tugas/praktik mandiri. & Mahasiswa mengerjakan aktivitas terarah untuk standar keberlanjutan TI dan menunjukkan evidence hasil belajar. & Bentuk penilaian: Standar Keberlanjutan TI. Mengacu rubrik RTM. & Ketepatan konsep; kualitas analisis; validasi dan dokumentasi. & 3 \\
4 & SCPMK0101-05201 & Kuis 1 konsep dan regulasi komputasi hijau. & Modalitas: Blended Learning. Bentuk: Luring/Daring. Metode: Case Method, studi kasus, diskusi lingkungan, dan peer review. & 1 x 2 x 50' tatap muka; 1 x 2 x 50' tugas/praktik mandiri. & Mahasiswa mengerjakan aktivitas terarah untuk kuis 1 konsep dan regulasi komputasi hijau dan menunjukkan evidence hasil belajar. & Bentuk penilaian: Kuis 1 Konsep dan Regulasi Komputasi Hijau. Mengacu rubrik RTM. & Ketepatan konsep; kualitas analisis; validasi dan dokumentasi. & 5 \\
5 & SCPMK0509-05202 & Efisiensi energi perangkat keras: server, pendingin, UPS. & Modalitas: Blended Learning. Bentuk: Luring/Daring. Metode: Case Method, studi kasus, diskusi lingkungan, dan peer review. & 1 x 2 x 50' tatap muka; 1 x 2 x 50' tugas/praktik mandiri. & Mahasiswa mengerjakan aktivitas terarah untuk efisiensi energi perangkat keras dan menunjukkan evidence hasil belajar. & Bentuk penilaian: Efisiensi Energi Perangkat Keras. Mengacu rubrik RTM. & Ketepatan konsep; kualitas analisis; validasi dan dokumentasi. & 2 \\
6 & SCPMK0509-05202 & Power Usage Effectiveness (PUE) dan data center hijau. & Modalitas: Blended Learning. Bentuk: Luring/Daring. Metode: Case Method, studi kasus, diskusi lingkungan, dan peer review. & 1 x 2 x 50' tatap muka; 1 x 2 x 50' tugas/praktik mandiri. & Mahasiswa mengerjakan aktivitas terarah untuk PUE dan data center hijau dan menunjukkan evidence hasil belajar. & Bentuk penilaian: PUE dan Data Center Hijau. Mengacu rubrik RTM. & Ketepatan konsep; kualitas analisis; validasi dan dokumentasi. & 3 \\
7 & SCPMK0509-05202 & Virtualisasi dan konsolidasi server untuk hemat energi. & Modalitas: Blended Learning. Bentuk: Luring/Daring. Metode: Case Method, studi kasus, diskusi lingkungan, dan peer review. & 1 x 2 x 50' tatap muka; 1 x 2 x 50' tugas/praktik mandiri. & Mahasiswa mengerjakan aktivitas terarah untuk virtualisasi dan konsolidasi server dan menunjukkan evidence hasil belajar. & Bentuk penilaian: Virtualisasi dan Konsolidasi Server. Mengacu rubrik RTM. & Ketepatan konsep; kualitas analisis; validasi dan dokumentasi. & 3 \\
8 & SCPMK0101-05201, SCPMK0509-05202 & UTS carbon footprint dan efisiensi energi. & Modalitas: Blended Learning. Bentuk: Luring/Daring. Metode: Case Method, studi kasus, diskusi lingkungan, dan peer review. & 1 x 2 x 50' tatap muka; 1 x 2 x 50' tugas/praktik mandiri. & Mahasiswa mengerjakan aktivitas terarah untuk UTS carbon footprint dan efisiensi energi dan menunjukkan evidence hasil belajar. & Bentuk penilaian: UTS Carbon Footprint dan Efisiensi Energi. Mengacu rubrik RTM. & Ketepatan konsep; kualitas analisis; validasi dan dokumentasi. & 15 \\
9 & SCPMK0509-05202 & Cloud computing sebagai enabler komputasi hijau. & Modalitas: Blended Learning. Bentuk: Luring/Daring. Metode: Case Method, studi kasus, diskusi lingkungan, dan peer review. & 1 x 2 x 50' tatap muka; 1 x 2 x 50' tugas/praktik mandiri. & Mahasiswa mengerjakan aktivitas terarah untuk cloud computing sebagai enabler komputasi hijau dan menunjukkan evidence hasil belajar. & Bentuk penilaian: Cloud Computing sebagai Enabler Komputasi Hijau. Mengacu rubrik RTM. & Ketepatan konsep; kualitas analisis; validasi dan dokumentasi. & 2 \\
10 & SCPMK0509-05202 & Green software engineering. & Modalitas: Blended Learning. Bentuk: Luring/Daring. Metode: Case Method, studi kasus, diskusi lingkungan, dan peer review. & 1 x 2 x 50' tatap muka; 1 x 2 x 50' tugas/praktik mandiri. & Mahasiswa mengerjakan aktivitas terarah untuk green software engineering dan menunjukkan evidence hasil belajar. & Bentuk penilaian: Green Software Engineering. Mengacu rubrik RTM. & Ketepatan konsep; kualitas analisis; validasi dan dokumentasi. & 3 \\
11 & SCPMK0509-05202 & Kuis 2 virtualisasi dan green software. & Modalitas: Blended Learning. Bentuk: Luring/Daring. Metode: Case Method, studi kasus, diskusi lingkungan, dan peer review. & 1 x 2 x 50' tatap muka; 1 x 2 x 50' tugas/praktik mandiri. & Mahasiswa mengerjakan aktivitas terarah untuk kuis 2 virtualisasi dan green software dan menunjukkan evidence hasil belajar. & Bentuk penilaian: Kuis 2 Virtualisasi dan Green Software. Mengacu rubrik RTM. & Ketepatan konsep; kualitas analisis; validasi dan dokumentasi. & 5 \\
12 & SCPMK0101-05201 & E-waste: daur ulang dan pengelolaan perangkat TI. & Modalitas: Blended Learning. Bentuk: Luring/Daring. Metode: Case Method, studi kasus, diskusi lingkungan, dan peer review. & 1 x 2 x 50' tatap muka; 1 x 2 x 50' tugas/praktik mandiri. & Mahasiswa mengerjakan aktivitas terarah untuk e-waste dan pengelolaan perangkat TI dan menunjukkan evidence hasil belajar. & Bentuk penilaian: E-waste dan Pengelolaan Perangkat TI. Mengacu rubrik RTM. & Ketepatan konsep; kualitas analisis; validasi dan dokumentasi. & 3 \\
13 & SCPMK0509-05202 & Renewable energy untuk infrastruktur TI. & Modalitas: Blended Learning. Bentuk: Luring/Daring. Metode: Case Method, studi kasus, diskusi lingkungan, dan peer review. & 1 x 2 x 50' tatap muka; 1 x 2 x 50' tugas/praktik mandiri. & Mahasiswa mengerjakan aktivitas terarah untuk renewable energy untuk infrastruktur TI dan menunjukkan evidence hasil belajar. & Bentuk penilaian: Renewable Energy untuk Infrastruktur TI. Mengacu rubrik RTM. & Ketepatan konsep; kualitas analisis; validasi dan dokumentasi. & 3 \\
14 & SCPMK0101-05201 & Green IT governance dan kebijakan organisasi. & Modalitas: Blended Learning. Bentuk: Luring/Daring. Metode: Case Method, studi kasus, diskusi lingkungan, dan peer review. & 1 x 2 x 50' tatap muka; 1 x 2 x 50' tugas/praktik mandiri. & Mahasiswa mengerjakan aktivitas terarah untuk green IT governance dan kebijakan organisasi dan menunjukkan evidence hasil belajar. & Bentuk penilaian: Green IT Governance dan Kebijakan Organisasi. Mengacu rubrik RTM. & Ketepatan konsep; kualitas analisis; validasi dan dokumentasi. & 2 \\
15 & SCPMK0509-05202 & Studi kasus: green data center dan sustainable IT. & Modalitas: Blended Learning. Bentuk: Luring/Daring. Metode: Case Method, studi kasus, diskusi lingkungan, dan peer review. & 1 x 2 x 50' tatap muka; 1 x 2 x 50' tugas/praktik mandiri. & Mahasiswa mengerjakan aktivitas terarah untuk studi kasus green data center dan menunjukkan evidence hasil belajar. & Bentuk penilaian: Studi Kasus Green Data Center. Mengacu rubrik RTM. & Ketepatan konsep; kualitas analisis; validasi dan dokumentasi. & 2 \\
16 & SCPMK0101-05201, SCPMK0509-05202 & PBL audit green IT dan rekomendasi. & Modalitas: Blended Learning. Bentuk: Luring/Daring. Metode: Case Method, studi kasus, diskusi lingkungan, dan peer review. & 1 x 2 x 50' tatap muka; 1 x 2 x 50' tugas/praktik mandiri. & Mahasiswa mengerjakan aktivitas terarah untuk PBL audit green IT dan rekomendasi dan menunjukkan evidence hasil belajar. & Bentuk penilaian: PBL Audit Green IT dan Rekomendasi. Mengacu rubrik RTM. & Ketepatan konsep; kualitas analisis; validasi dan dokumentasi. & 22 \\
17 & SCPMK0509-05202 & UAS presentasi desain infrastruktur berkelanjutan. & Modalitas: Blended Learning. Bentuk: Luring/Daring. Metode: Case Method, studi kasus, diskusi lingkungan, dan peer review. & 1 x 2 x 50' tatap muka; 1 x 2 x 50' tugas/praktik mandiri. & Mahasiswa mengerjakan aktivitas terarah untuk UAS presentasi desain infrastruktur berkelanjutan dan menunjukkan evidence hasil belajar. & Bentuk penilaian: UAS Presentasi Desain Infrastruktur Berkelanjutan. Mengacu rubrik RTM. & Ketepatan konsep; kualitas analisis; validasi dan dokumentasi. & 23 \\
\end{longtblr}
